\scanpage{124}
de réalisabilité nous donne en $\overline{\mathrm m}$ temps que $A\simeq M(b,\unsure{c|b},b',c',t)$, laquelle est donc bien Azumaya. On gagne !

Mais j'ai oublié de faire la vérification (sur le \deleted{alg} corps alg clos) quand $\delta\delta'=0$, donc $H$ ou $H'$ est ``unipotente'' \deleted{\unsure{i.e. on}} (disons $H'$) i.e. \deleted{Considérons le discriminant de} \deleted{la forme} (comme on est sur $k$ alg clos) \deleted{$Q(t)$ considéré comme un polynôme en} $H'=\deleted{\unsure{k[v]/v^2}}k[V]/V^2$, $V$ une indét. \deleted{$L\subset A^{\unsure{\mathrm{s}|\mathrm{tn}}}$ tot isotr. i.e. $\det|L=0$ $\operatorname{Tr}|L=0$}

La formulation qui suit montre que pour $H,H'$ donnés, la \inserted{réponse au} problème de réalisabilité ne dépend pas du choix des \deleted{bases} éléments $u\in H$, $v\in H'$ utilisés pour faire des bases. \deleted{On peut donc} \inserted{Il y a donc} trois ``cas types''

1) $\delta$ \uline{et} $\delta'$ \deleted{\unsure{$\ne1$}}$\ne0$ \quad OPS $u^2=u,v^2=v$ (idempotents : donc $b=b'=1$, $c=c'=0$
\[
Q(t)=t^2-t=(t-1)t
\]

2) $\delta\deleted{\unsure{\text{ et }\delta'}}\ne0,\ \delta'\deleted{\unsure{\ne0}}=0$ (ou l'inverse) OPS \quad $u^2=u\quad v^2=\unsure{0|v}$ i.e. $\begin{aligned}b&=1&c&=0\\b'&=0&c'&=0\end{aligned}$

3) $\delta=\delta'=0$ \quad OPS \quad $u^2=0,v^2=0$ i.e $b=b'=c=c'=0$

Dans les cas $2^\circ)$ et $3^\circ)$ et dès qu'on prend $v^2=0$, (sans condition sur $u$) on trouve
\[
Q(t)=t^2
\]
% archive sheet number: 123
